\documentclass[11pt]{article}
\usepackage{amsmath}
\usepackage{amsfonts}
\usepackage{amsthm}
\usepackage[utf8]{inputenc}
\usepackage[margin=0.75in]{geometry}
\usepackage{hyperref}

\title{CSC111 Winter 2026 Project 2:\\Charting Your Academic Path:
A Graph-Based Course Planning Tool for UofT Students}
\author{Jack Anderson, Efren Medina Arias, Diego Jose Figarella Urbaneja, Tanish Ariyur}
\date{\today}

\begin{document}
\maketitle

\section*{Introduction}

The University of Toronto offers hundreds of courses across the Faculty of Arts and Science,
many of which sit behind long chains of prerequisite requirements. As students at UofT,
planning our paths toward our degrees is often tedious and confusing. A student who wants to
eventually enrol in an upper-year course has to work backwards through multiple layers of
prerequisites, often discovering gaps in their plan only after the fact. This adds unnecessary
stress and uncertainty to academic planning, regardless of what program or discipline a student
is pursuing.

This problem maps naturally onto a directed graph, where each course is a vertex and each
prerequisite relationship is a directed edge from the required course to the one it unlocks.
UofT prerequisites are not simple enough to be stored as plain directed edges, though, since
a prerequisite string like ``(CSC148H1 AND CSC165H1) OR CSC240H1'' is a logical condition,
not a single course. Our implementation reflects this: \verb|CourseGraph| stores courses as
vertices, but instead of explicit edges, each vertex holds a \verb|BooleanList| prerequisite
tree that encodes which other courses unlock it. These \verb|BooleanList| trees are the tree
structure at the core of the project. The course codes appearing as leaves of a
\verb|BooleanList| implicitly define directed relationships between courses (what we call
pseudo-edges), since each leaf represents a course that must be completed before the current
one. The full graph of prerequisite-of relationships is therefore encoded in the collection of
these trees, one per vertex, rather than in an explicit adjacency structure. When needed for
visualization, \verb|CourseGraph.to_networkx()| materializes these pseudo-edges into real
directed edges by extracting all course codes from each vertex's \verb|BooleanList| via
\verb|get_all_courses()| and adding an edge from each to the current course.

\textbf{Given a student's completed courses and a target course they wish to take, our
program determines the prerequisite courses the student needs to complete to satisfy all
requirements, guides them through selecting those courses interactively, and displays the
prerequisite network in an interactive visual format so the student can see the full picture
of what they are working toward.}

This problem has been studied in the context of academic advising systems, where course
catalogues are typically modelled as directed acyclic graphs [2]. Our approach is consistent
with that framing, with the added complexity of UofT's prerequisite strings, which mix AND
and OR conditions and credit thresholds in a single natural-language string.

\section*{Dataset Description}

Our primary data source is the University of Toronto Faculty of Arts and Science 2025--2026
Academic Calendar [1]. We worked with an HTML version of the calendar, which contains for
each course its course code, title, description, contact hours, breadth requirement,
prerequisite requirements, and exclusion notices. A representative example from the calendar
is:

\begin{verbatim}
CSC236H1 - Introduction to the Theory of Computation.
Prerequisite: (60% or higher in CSC111H1) and (60% or higher in
             MAT102H5 or MAT137Y1 or MAT157Y1)
\end{verbatim}

We used the Python library \verb|beautifulsoup4| [5] to parse the HTML calendar file and
extract each course entry into a structured JSON file. Concretely, \verb|parse.py| uses
BeautifulSoup's \verb|find| method to locate course entries within \verb|class='view-content'|,
then extracts the title, body, breadth, hours, prerequisite, and exclusion fields from each
entry's sub-elements. The code and name are split from the title string, and the breadth
number is parsed from the trailing parenthetical in that field. The result is written to
\verb|data/raw/courses.json|, a clean, reusable dataset that we can load without
re-parsing the HTML each time.

The harder problem was converting the raw prerequisite strings into something the program can
actually evaluate. A string like \verb|"CSC148H1 and CSC165H1/ CSC240H1"| has to become a
structured AND/OR tree. The more complex prerequisite strings (those with deep nesting,
mixed conjunctions, or credit-threshold clauses) were post-processed using large language
models (Claude and ChatGPT), an approach we received explicit approval for from our TA and
professor. The final processed file used by the program at runtime is
\verb|data/courses_formatted_3.json|, where each entry's \verb|prereq_tree| field encodes
the full logical structure of its prerequisites as nested AND/OR nodes and credit-threshold
nodes.

At runtime, \verb|json_to_graph.py| reads this file and loads each entry's course code,
name, hours, description, breadth number, \verb|prereq_tree|, and exclusions into the graph.
The raw prerequisite string is not kept; only the structured tree.

\section*{Computational Overview}

\textbf{How graphs and trees play a central role.} The entire program is built on top of
\verb|CourseGraph| in \verb|course_graph.py|, a directed graph where each vertex is a
\verb|_CourseVertex| storing all course metadata. Every core operation (checking whether a
student can take a course, finding what they should take next, enumerating all paths to a
target, counting credits) is a query on this graph. There is no way to run the program
without constructing and traversing it.

The connections between vertices are encoded as \verb|BooleanList| prerequisite trees, one
per vertex, defined in \verb|boolean_list.py|. \verb|BooleanList| is a fully custom-built
recursive data structure, not backed by any external library, designed specifically to
represent and evaluate the propositional logic of UofT prerequisite conditions. It is the
core innovation of the project. Without it, none of the eligibility checking, path-finding,
or interactive navigation is possible, because UofT's prerequisites cannot be reduced to a
flat list of courses.

\textbf{Structure.} Each \verb|BooleanList| node has an \verb|operator| (either \verb|'AND'|
or \verb|'OR'|) and an \verb|items| list whose elements can be one of three types:
\begin{itemize}
    \item A \textbf{course code string}: a leaf representing a single course that must
    be completed. Each such string is also a pseudo-edge in the graph: it encodes a directed
    prerequisite-of relationship from that course to the vertex holding the tree, without
    requiring an explicit entry in an adjacency structure.
    \item A \textbf{nested} \verb|BooleanList|: a sub-condition with its own operator and
    items, allowing arbitrarily deep nesting. This is what lets the structure represent
    conditions like ``(CSC148H1 AND CSC165H1) OR CSC240H1'' faithfully: the outer node is an
    OR, and one of its children is an AND node containing the two courses.
    \item A \textbf{\texttt{CreditCondition}} node: a custom class representing a
    quantitative credit-threshold requirement such as ``at least 1.0 credit in CSC.''
    \verb|CreditCondition| is also fully custom: \verb|credits_satisfied()| infers credit
    weight directly from the course code suffix (\verb|H| = 0.5, \verb|Y| = 1.0) with no
    external lookup, and \verb|calculate_credits_department()| filters by department prefix to
    handle requirements like ``1.0 FCE in MAT.''
\end{itemize}

\textbf{Evaluation.} \verb|BooleanList.is_satisfied(completed)| is a recursive propositional
logic evaluator: given a set of completed course codes, it traverses the tree and returns
whether the entire formula is true. The dispatch to each item type is handled by the static
method \verb|evaluate_item(item, completed)|, which branches on \verb|isinstance| checks.
A string item is satisfied if it is in \verb|completed|, a \verb|CreditCondition| item
delegates to \verb|credits_satisfied()|, and a nested \verb|BooleanList| item recurses
back into \verb|is_satisfied()|. AND nodes use Python's \verb|all()| over their items, OR
nodes use \verb|any()|, both of which short-circuit. The result is a full propositional model
checker for UofT prerequisite formulas, written from scratch using only the language features
covered in CSC111.

This design is what makes a condition like the CSC236H1 prerequisite,
\textit{``(60\% or higher in CSC111H1) and (60\% or higher in MAT102H5 or MAT137Y1 or MAT157Y1)''},
evaluable in a single call to \verb|is_satisfied()|. The outer AND node has two children:
the first is a leaf for \verb|CSC111H1|, and the second is an OR node containing
\verb|MAT102H5|, \verb|MAT137Y1|, and \verb|MAT157Y1|. No other standard data structure
from the course (linked lists, plain trees, or flat dicts) could represent this without
losing information about which combination of courses satisfies the requirement.

\textbf{Pseudo-edge extraction.} \verb|BooleanList.get_all_courses()| recursively collects
every course code string from any depth of the tree, skipping \verb|CreditCondition| nodes
since they do not name a specific course. This is used by \verb|CourseGraph.to_networkx()|,
which calls it on each vertex's \verb|BooleanList| and adds a directed edge from every
returned code to the current course, materializing the pseudo-edges into real edges in a
\verb|nx.DiGraph| for visualization. Outside of \verb|to_networkx()|, no explicit edge set
is ever constructed; all prerequisite relationships are accessed purely by evaluating the
\verb|BooleanList| trees in place.

\textbf{Major computations.}

\begin{enumerate}

    \item \textit{Eligibility checking.}
    \verb|CourseGraph.is_eligible(code, completed)| returns whether a student who has
    completed the given set of courses can enrol in \verb|code|. It returns \verb|False|
    immediately if the student has already completed the course, \verb|True| if the course
    has no prerequisites, and otherwise delegates to
    \verb|BooleanList.is_satisfied(completed)|, which recursively evaluates each item via
    \verb|evaluate_item()|. \verb|eligible_courses(completed)| calls this across every vertex
    to return the full set of currently available courses.

    \item \textit{Next-course resolution.}
    \verb|get_next_needed_courses(graph, target, completed)| in \verb|algorithms.py| is the
    algorithm that drives the interactive interface. Given a target course the student cannot
    yet take, it walks the prerequisite tree to find the courses they should pick from right
    now to make progress. If the target is already eligible, it returns the target itself. If
    not, it passes the target's \verb|BooleanList| to \verb|_resolve_needed()|, which
    dispatches to \verb|_resolve_and()| or \verb|_resolve_or()| depending on the node type.

    For AND nodes, \verb|_resolve_and()| collects suggestions from all unsatisfied children
    simultaneously, since the student has to complete all of them anyway. If a child course
    is not yet eligible itself, the method recurses into its own prerequisite tree to find
    what needs to happen first.

    OR nodes are the tricky part. When a prerequisite is ``CSC165H1 or CSC240H1,'' the
    student who has already started taking CSC165H1-level courses should not suddenly be
    shown CSC240H1 as a new option. \verb|_resolve_or()| handles this by first checking
    whether any branch is already partially satisfied via \verb|_has_progress()|, which
    returns \verb|True| if at least one sub-item in a branch is in \verb|completed|. If
    there are branches with prior progress, \verb|_partition_or_children()| separates them
    from the rest and \verb|_resolve_or()| focuses only on those branches, ignoring the
    others entirely. This keeps the suggestions relevant rather than flooding the student
    with every possible alternative every time they take a course.

    A \verb|visited| set is threaded through all recursive calls to prevent infinite loops
    when the prerequisite graph has cycles.

    \item \textit{Relevance filtering.}
    \verb|get_relevant_courses(graph, target, completed)| computes the set of all courses
    that appear on any prerequisite path to \verb|target|, stopping recursion at courses
    already in \verb|completed|. \verb|eligible_relevant_courses()| intersects this with the
    eligible courses to return only courses that are both takeable now and actually useful for
    reaching the target.

    \item \textit{Credit counting.}
    \verb|CourseGraph.credit_count(completed, department)| returns the total credits
    accumulated from a set of completed courses, with an optional filter by department prefix.
    Credit weight is inferred from the course code suffix (\verb|H| = 0.5, \verb|Y| = 1.0).
    This is useful for checking whether program-level credit requirements are met.

    \item \textit{Course search.}
    \verb|search_courses(graph, query)| returns courses whose code or name contains the query
    string (case-insensitive), sorted so that exact code prefix matches come first. This backs
    the live search bars in the interface.

\end{enumerate}

\textbf{Result reporting.} The program presents results through two interfaces.

The main interface is \verb|CourseNavigator| in \verb|interface.py|, a Tkinter application
with a setup phase and a navigation phase. In the setup phase the student adds completed
courses via a live search bar and types a target course. Clicking ``Start navigating''
switches to the navigation phase. The centre panel shows course cards generated by
\verb|get_next_needed_courses()|, sorted by department and level, with the target course
highlighted in yellow if it is already reachable. Clicking a card adds that course to
\verb|completed| and \verb|path|, then refreshes the options. The left panel tracks the path
taken so far with a step counter and an undo button. The right panel shows course metadata
on hover: credit weight, level, department, breadth category, a formatted prerequisite
tree with checkmarks beside completed items, the exclusions list, and a clickable link to the
course's page on the UofT calendar [1]. The session ends on a success screen when the target
is added to \verb|completed|, showing the full path taken and the total number of courses
selected. \verb|CourseNavigator| carries 18 instance attributes in total, one for each
persistently referenced widget, so that every attribute is fully type-annotated in the
class body. This exceeds \verb|python_ta|'s default limit of 7; the \verb|max-attributes|
threshold was raised to 20 in the \verb|python_ta| configuration for \verb|interface.py|
to accommodate the complete set of type annotations.

The visualization is \verb|visualize_course_graph()| in \verb|visualizations.py|. It renders
the full course prerequisite graph as an interactive Plotly [3] figure. Node positions are
computed by a custom hierarchical layout: \verb|get_depths()| assigns each course a depth
equal to the length of the longest prerequisite chain leading to it via topological sort
(using NetworkX [6]), and \verb|hierarchical_layout()| places courses at the same depth on
the same vertical level. Two \verb|go.Scatter()| traces are combined, one with
\verb|mode='lines'| for edges and one with \verb|mode='markers+text'| for nodes, and
rendered via \verb|fig.show()|.

\textbf{New Python libraries.}

\begin{itemize}
    \item \verb|networkx| [6]: converts \verb|CourseGraph| to a \verb|nx.DiGraph| via
    \verb|CourseGraph.to_networkx()|; used in \verb|visualizations.py| for topological
    sorting and hierarchical depth assignment.
    \item \verb|plotly| [3]: renders the interactive course graph in
    \verb|visualizations.py| using \verb|go.Scatter()| traces and \verb|go.Figure()|.
    \item \verb|beautifulsoup4| [5]: parses the HTML course calendar in \verb|parse.py|
    to extract structured course data into JSON.
\end{itemize}

\section*{Running Instructions}

\begin{enumerate}
    \item Install Python 3.13.5.
    \item Install dependencies:
\begin{verbatim}
pip install beautifulsoup4 lxml networkx plotly
\end{verbatim}
    \item Run the program:
\begin{verbatim}
python main.py
\end{verbatim}
\end{enumerate}

The processed dataset (\verb|data/courses_formatted_3.json|) is included with the submission,
so no additional data download or preprocessing is required. The Tkinter window opens directly
on launch.

To view the graph visualization instead, run \verb|visualizations.py| directly
(\verb|python visualizations.py|), which opens an interactive Plotly figure in the browser.

\textbf{Example walkthrough: reaching CSC311H1 (Introduction to Machine Learning).}
CSC311H1 requires, via the CS pathway: a programming background through CSC110Y1 and
CSC111H1, then CSC207H1, plus calculus (MAT137Y1), linear algebra (MAT223H1), and a
probability course (STA237H1). In the setup phase, leave the completed courses list empty,
type \verb|CSC311H1| in the target field, and click ``Start navigating.''

The navigation phase presents the first layer of actionable courses:
\verb|CSC110Y1|, \verb|MAT137Y1|, and \verb|MAT223H1| (all have no prerequisites and are all
needed). Click \verb|CSC110Y1|. The options update to include \verb|CSC111H1| (now eligible),
\verb|MAT137Y1|, and \verb|MAT223H1|. Continue clicking through the suggested cards in order:
\verb|MAT137Y1|, \verb|MAT223H1|, \verb|CSC111H1|, \verb|STA237H1| (eligible once
\verb|MAT137Y1| is done), \verb|CSC207H1|, and finally \verb|CSC311H1| itself, which appears
highlighted in yellow as the target once all prerequisites are satisfied. The success screen
displays the full 7-course path:
\begin{verbatim}
CSC110Y1 -> MAT137Y1 -> MAT223H1 -> CSC111H1 -> STA237H1 -> CSC207H1 -> CSC311H1
\end{verbatim}

\section*{Changes to Project Plan}

\begin{itemize}

    \item \textbf{Graph and tree models both used.} The proposal described both a directed
    graph and a \verb|CourseTree| and noted we were still figuring out which to use. The
    final implementation uses both. \verb|CourseGraph| is the graph, with courses as vertices.
    The \verb|BooleanList| on each vertex is the tree: a recursive AND/OR tree that encodes
    the prerequisite logic for that course and implicitly defines directed relationships
    (pseudo-edges) to the courses it lists as leaves. A separate \verb|CourseTree| class as
    described in the proposal was not implemented; the tree structure lives inside the
    graph rather than alongside it. The student state (completed courses) is passed as a
    parameter to graph methods rather than being embedded in tree nodes, which avoids the
    node-explosion problem the proposal identified.

    \item \textbf{BooleanList and CreditCondition added.} The proposal described extracting
    course codes from prerequisite strings via regex and storing ``multiple valid prerequisite
    combinations per vertex.'' This turned out to be insufficient. A flat list of combinations
    cannot distinguish ``CSC148H1 AND (CSC165H1 OR CSC240H1)'' from
    ``(CSC148H1 AND CSC165H1) OR CSC240H1,'' since the grouping information is lost. The
    final implementation replaces this with \verb|BooleanList|, a fully custom recursive
    AND/OR tree that preserves the full logical structure of each prerequisite condition and
    evaluates it directly against the student's completed set. \verb|CreditCondition| was
    added alongside it to handle the subset of requirements expressed as credit thresholds
    rather than specific course codes. Neither class exists in any library; both are written
    from scratch specifically for this project.

    \item \textbf{Corequisites not implemented.} The proposal mentioned storing corequisites
    as a distinct edge type on the graph. In practice, UofT's calendar does not consistently
    separate corequisites from prerequisites in the HTML, so we did not implement a distinct
    corequisite edge type. Corequisite requirements that appeared in the prerequisite field
    were treated as regular prerequisites during parsing.

    \item \textbf{Command-line input replaced by Tkinter.} The proposal described a
    command-line interface for entering completed courses and target courses. The final
    implementation uses a Tkinter GUI with live search bars and clickable course cards, which
    is more usable and avoids requiring students to type exact course codes from memory.

    \item \textbf{Visualization colour scheme simplified.} The proposal planned to
    colour-code nodes as green (completed), yellow (on the recommended path), and red
    (target). The final visualization renders all nodes with the same colour since the
    full-graph view is not filtered to a student's specific path. The per-student
    colour-coding was moved into the Tkinter interface instead, where the target course card
    is highlighted in yellow and path items in the left panel are coloured green.

    \item \textbf{find\_path\_to $\rightarrow$ get\_next\_needed\_courses.} The proposal
    described a \verb|find_path_to(target)| method that returns all valid sequences to a
    target. The final implementation replaces this with \verb|get_next_needed_courses()| in
    \verb|algorithms.py|, which resolves the courses the student should pick from at each
    step interactively rather than computing all paths upfront. This fits the layer-by-layer
    navigation model of the interface better than a full path enumeration would.

    \item \textbf{get\_credits $\rightarrow$ credit\_count.} \verb|get_credits(optional[Dept])|
    from the proposal became \verb|CourseGraph.credit_count(completed, department)| with the
    same semantics.

\end{itemize}

\section*{Discussion}

The problem we set out to solve was one we ran into personally as UofT students, figuring
out what courses you need to take before you can take the one you actually want. Working
backwards through the calendar manually is tedious enough when prerequisites are simple, but
UofT's prerequisite strings are often nested combinations of AND and OR conditions, sometimes
with credit thresholds mixed in. The goal was to make that process something the program
handles instead of the student.

The single most important design decision in the project was building \verb|BooleanList|.
Every other component (the graph, the navigation algorithm, the interface) depends on
it. The insight behind it is that UofT prerequisite conditions are propositional logic
formulas over course codes, and the right way to represent a logic formula is as a tree.
Standard data structures from the course (linked lists, plain trees, flat dicts) cannot
express ``complete A AND B, OR just C'' without collapsing the OR and the AND into something
ambiguous. \verb|BooleanList| solves this by being a fully recursive AND/OR tree where each
node carries an operator and a list of items that can themselves be \verb|BooleanList| nodes,
giving it the ability to represent arbitrary nesting depth.

What makes \verb|BooleanList| non-trivial to build is the three-way polymorphism of its
items. A node's children can be course code strings, nested \verb|BooleanList| nodes, or
\verb|CreditCondition| nodes, all handled uniformly by \verb|evaluate_item()|, which
dispatches on type. This means \verb|is_satisfied()| never needs to know in advance what
kind of item it is evaluating; it just calls \verb|evaluate_item()| on each child and lets
the dispatch handle the rest. Getting this right required thinking carefully about when
recursion terminates (at string leaves and \verb|CreditCondition| nodes, which do not recurse
further) and making sure AND/OR short-circuiting via \verb|all()| and \verb|any()| did not
skip evaluations that had side effects elsewhere in the tree.

\verb|CreditCondition| is its own custom class rather than a special string format, which
matters because credit requirements need to be evaluated against the student's accumulated
credits, not against a set of course codes. \verb|credits_satisfied()| infers whether a
student meets the threshold by parsing credit weight directly from course code suffixes
(\verb|H| for 0.5, \verb|Y| for 1.0), with an optional department filter, keeping the
entire eligibility check self-contained.

The hardest part of getting \verb|BooleanList| into the program was not writing the class
itself but producing valid trees from the raw calendar data. The prerequisite strings from
the HTML are formatted inconsistently: some use commas for AND, some use semicolons, some mix
``or'' and ``/'' in the same clause. The more complex strings (those with deep nesting,
mixed conjunctions, and credit-threshold clauses in the same expression) were
post-processed using large language models (Claude and ChatGPT) to produce valid
\verb|prereq_tree| dictionaries. We received explicit approval from our TA and professor to
use this approach for the data processing step. Without a validated conversion, malformed
trees would silently enter the graph and cause incorrect eligibility results across the entire
dataset.

The \verb|get_next_needed_courses()| algorithm in \verb|algorithms.py| was the most
interesting piece of code to write, and it is built entirely on top of \verb|BooleanList|.
It works by recursing through the prerequisite tree of the target course and asking, at each
node, what does the student need to take right now to make progress? The AND case is
straightforward (collect from all unsatisfied children at once), but OR nodes required
more thought. The issue is that once a student has started taking courses down one branch of
an OR prerequisite, you do not want to suggest the other branches anymore. If someone has
already taken \verb|CSC165H1| as part of getting to \verb|CSC207H1|, showing them
\verb|CSC240H1| as an alternative at that point is not useful. \verb|_has_progress()| and
\verb|_partition_or_children()| solve this by detecting whether any sub-item in a branch is
already satisfied and, if so, restricting suggestions to those branches only. Getting this
right took a few iterations; early versions would open up all branches simultaneously and
produce suggestions that had nothing to do with where the student was in their path. The fix
depended on being able to introspect the \verb|BooleanList| tree structurally, asking not
just ``is this satisfied?'' but ``does any sub-item show progress?'', which is only
possible because the tree is a custom data structure we built and control entirely.

One thing we did not get to is integrating exclusions into eligibility checking. The exclusion
list is stored on every \verb|_CourseVertex| and is displayed in the right panel of the
interface, but \verb|is_eligible()| currently only checks prerequisites, not exclusions. This
means the program could in theory suggest a course that a student is technically excluded from
because they took an equivalent. Fixing this is straightforward: \verb|is_eligible()| would
just need to return \verb|False| if any entry in \verb|vertex.exclusions| appears in
\verb|completed|, but we ran out of time to add it cleanly given the number of edge cases
in how the exclusion strings are formatted in the raw data.

The Plotly visualization works well for smaller subsets of the graph, for example all
CSC-prefixed courses or the prerequisite chain for a single target. At the full scale of
roughly 1,300 courses it becomes too dense to read. A natural extension would be to filter the
visualization to only the courses returned by \verb|get_relevant_courses()|, which already
computes exactly the subgraph relevant to a given target. Passing that directly to
\verb|visualize_course_graph()| would give the student a readable view of just the courses
that matter for their goal, coloured by whether they have been completed, which was the
original intent of the colour-coding scheme from the proposal.

Overall, the core of what we set out to build works: a student can open the program, type in
a target like \verb|CSC373H1|, add the courses they have already taken, and then click
through the suggested course cards until the target appears as an option. The prerequisite
tree panel in the right sidebar shows exactly where they are in the requirement chain, with
checkmarks beside the parts they have already satisfied. That was the original goal, and it
does that reliably.

\section*{References}

\begin{enumerate}

    \item[{[1]}] University of Toronto Faculty of Arts \& Science. ``2025--2026 Academic
    Calendar.'' \textit{University of Toronto}, 2025.
    Available at: \url{https://artsci.calendar.utoronto.ca/}. Accessed March 2026.

    \item[{[2]}] Barachini, F., and Stary, C. ``From Knowledge Representation to Knowledge
    Communication: Course Prerequisite Graphs in Academic Advising.'' \textit{Proceedings of
    the 2nd International Conference on Computer Supported Education}, 2010, pp.\ 135--144.

    \item[{[3]}] Plotly Technologies Inc. ``Plotly Python Open Source Graphing Library.''
    \textit{Plotly}, 2025. Available at: \url{https://plotly.com/python/}. Accessed March
    2026.

    \item[{[4]}] Plotly Technologies Inc. ``Network Graphs in Python.'' \textit{Plotly},
    2025. Available at: \url{https://plotly.com/python/network-graphs/}. Accessed March 2026.

    \item[{[5]}] Richardson, L. ``Beautiful Soup Documentation.'' \textit{Crummy}, 2025.
    Available at: \url{https://www.crummy.com/software/BeautifulSoup/bs4/doc/}. Accessed
    March 2026.

    \item[{[6]}] Hagberg, A., Swart, P., and Chult, D.\ S. ``Exploring Network Structure,
    Dynamics, and Function using NetworkX.'' \textit{Proceedings of the 7th Python in Science
    Conference}, 2008, pp.\ 11--15.

\end{enumerate}

\end{document}
